% Part: proof-theory
% Chapter: natural-deduction
% Section: rules-proofs

\documentclass[../../../include/open-logic-section]{subfiles}

\begin{document}

\olfileid{pt}{nat}{rp}

\olsection{القواعد و\usetoken{P}{proof}}

نبدأ ببعض التعريفات وتثبيت بعض المصطلحات.

تأمل القاعدتين في~\Log{N1c} و\Log{N1i} اللتين تذكران
$!A \lif !B$:
\[
\bottomAlignProof
\AxiomC{$\Discharge{!A}{x}$}
\shortDeduce
\DeduceC{$!B$}
\DischargeRule{\Intro{\lif}}{x}
\UnaryInfC{$!A \lif !B$}
\DisplayProof
\qquad
\bottomAlignProof
\AxiomC{$!A \lif !B$}
\AxiomC{$!A$}
\RightLabel{\Elim{\lif}}
\BinaryInfC{$!B$}
\DisplayProof
\]
قاعدة~$\Elim{\lif}$ بسيطة بما يكفي. تسمى~!!{formula}s الواقعة فوق
الخط \emph{المقدمات}. وتكون~!!{formula} الواقعة في اليسار هي
!!{formula}~$!A \lif !B$ \emph{الرئيسة}. ولكل قواعد~!!{elimination}
مقدمة من هذا النوع، تقع دائمًا في أقصى اليسار، وتسمى \emph{المقدمة
الكبرى} للاستدلال. وإذا وُجدت مقدمات أخرى، كما في هذه الحالة، سميت
\emph{مقدمات صغرى}. وتسمى مقدمات قواعد~!!{introduction} كذلك مقدمات
كبرى. أما~!!{formula} الواقعة تحت الخط فتسمى \emph{النتيجة}.

لقاعدة~$\Intro{\lif}$ مقدمة واحدة فقط، والنتيجة هنا هي الصيغة الرئيسة
$!A \lif !B$. وكل قواعد~!!{introduction} على هذا النحو: تكون النتيجة
هي الصيغة الرئيسة، ويكون عاملها الرئيس هو العامل الذي يجري «إدخاله».

يعني الترميز~$\Discharge{!A}{x}$ ما يأتي. في~!!a{proof}، يُطبق استدلال
$\Intro{\lif}$ على~!!a{formula}~$!B$ بوصفها مقدمة. وللـ!!{proof}
الجزئي المنتهي بـ$!B$ عدد من~!!{formula}s الابتدائية، وتسمى
\emph{فروضًا}. ويكون~!!^a{proof} لـ$!B$ من فرض أو أكثر على الصورة
$!A$~!!a{proof} على أن~$!A$ تستلزم~$!B$. وعندما نطبق قاعدة
$\Intro{\lif}$ نستنتج~$!A \lif !B$، ولا تعود النتيجة معتمدة على الفرض
$!A$. وللدلالة على ذلك، تُـ\emph{برأ} أو تُـ\emph{غلق} الفروض من
الصورة~$!A$ في~!!{proof} الذي يتضمن استدلال~$\Intro{\lif}$. وتحدد
القاعدة أي الفروض يمكن تبرئتها. ففي حالة~$\Intro{\lif}$ ذات النتيجة
$!A \lif !B$ تكون تلك الفروض على الصورة~$!A$، أي توافق مقدّم الشرط.
وقد يكون من المهم بيان أي الفروض تُبرأ عند أي قاعدة، ويؤدي علامة
التبرئة~$x$ هذا الغرض. ومن المهم أن القواعد المبرئة للفروض لا
\emph{يلزمها} تبرئة \emph{كل} الفروض ذات الصورة المطلوبة، بل لا يلزمها
تبرئة أي فرض أصلًا. فالآتي صحيح مثلًا، مع أن~$\Intro{\lif}$ الأولى لا
تبرئ أي فرض:
\[
\AxiomC{$\Discharge{!A}{y}$}
\DischargeRule{\Intro{\lif}}{x}
\UnaryInfC{$!B \lif !A$}
\DischargeRule{\Intro{\lif}}{y}
\UnaryInfC{$!A \lif (!B \lif !A)$}
\DisplayProof
\]
وعندما لا تبرئ قاعدة أي فرض نحذف علامة التبرئة ببساطة، كما فعلنا هنا.

في~!!{proof}
\[
\AxiomC{$\Discharge{!A}{x}$}
\AxiomC{$\Discharge{!A}{y}$}
\RightLabel{\Intro{\land}}
\BinaryInfC{$!A \land !A$}
\RightLabel{\Elim{\land}}
\UnaryInfC{$!A$}
\DischargeRule{\Intro{\lif}}{x}
\UnaryInfC{$!A \lif !A$}
\DischargeRule{\Intro{\lif}}{y}
\UnaryInfC{$!A \lif (!A \lif !A)$}
\DisplayProof
\]
يبرئ استدلال~$\Intro{\lif}$ الأول الفرض الأيسر~$!A^x$، ويبرئ الثاني
الفرض الأيمن~$!A^y$. أي إن كل الفروض الموسومة~$x$ يبرئها
$\Intro{\lif}$ ذو العلامة~$x$، وكل الفروض الموسومة~$y$ يبرئها
الاستدلال الموسوم~$y$. ولكي ينجح ذلك، من المهم أن تكون كل الفروض ذات
العلامة الواحدة~!!{formula} نفسها أيضًا.

قواعد الأسوار أعقد قليلًا. فقواعد~$\lexists$ في~\Log{N1i} مثلًا هي:
\[
\bottomAlignProof
\AxiomC{$!A(t)$}
\RightLabel{\Intro{\lexists}}
\UnaryInfC{$\lexists[x][!A(x)]$}
\DisplayProof
\qquad
\bottomAlignProof
\AxiomC{$\lexists[x][!A(x)]$}
\AxiomC{$\Discharge{!A(c)}{x}$}
\shortDeduce
\DeduceC{$!B$}
\DischargeRule{\Elim{\lexists}}{x}
\BinaryInfC{$!B$} \DisplayProof
\bottomAlignProof
\]
في قاعدة~$\Intro{\lexists}$ تكون المقدمة~$!A(t)$، حيث~$t$ حد مغلق،
أي حد بلا متغيرات. ويكون هذا الاستدلال صحيحًا، أي مطابقًا لمخطط قاعدة
$\Intro{\lexists}$، إذا وجدت~!!{formula}~$!A$ ومتغير~$x$ وحد مغلق
$t$ بحيث تكون النتيجة~$\lexists[x][!A]$ والمقدمة~$\Subst{!A}{t}{x}$.
وتتيح~$\Elim{\lexists}$ تبرئة فروض من الصورة~$!A(c)$؛ أي يوجد لكل
استدلال~!!{constant}~$c$ ويمكن تبرئة الفروض ذات الصورة
$\Subst{!A}{c}{x}$. ويجب أن تستخدم كل الفروض التي يبرئها استدلال واحد
الثابت~$c$ نفسه. ويسمى هذا الثابت \emph{المتغير المميز} لاستدلال
$\Elim{\lexists}$. ولا يصح الاستدلال إلا إذا لم يظهر~$c$ في فرض مفتوح
غير~$!A(c)$ المبرأ، ولا في~$!B$، ولا في~$!A$. ويسمى ذلك \emph{شرط
المتغير المميز}.

!!^{proof}s في~\Log{N1c} و\Log{N1i} أشجار منتهية من وقائع الصيغ، تُولّد
استقرائيًا من الفروض بتطبيق قواعد الاستدلال. وكل ورقة فرض، قد تكون
موسومة بعلامة تبرئة، وكل واقعة~!!{formula} أخرى تلزم من~!!{formula}s
الواقعة فوقها مباشرة بإحدى قواعد استدلال النظام. وإذا لم تكن واقعة
!!{formula} بوصفها فرضًا في الشجرة فوق استدلال قد بُرئت باستدلال فوقها،
سميت \emph{مفتوحة} عند ذلك الاستدلال.

\begin{defn}\ollabel{defn:proof}
تُعرّف \emph{!!^{proof}s} في~\Log{N1c} و\Log{N1i} استقرائيًا بوصفها
أشجارًا منتهية من~!!{formula}s كما يأتي:
\begin{enumerate}
  \item\ollabel{defn:prf:assumption} كل~!!{formula} موسومة، بمفردها،
  بعلامة تبرئة مثل~$x$ أو~$y$ هي~!!a{proof}. وفي~!!a{proof} يتألف من
  !!{formula} واحدة، تُعد~!!{formula} فرضًا مفتوحًا.
  \item\ollabel{defn:prf:one-prem} إذا كانت~$R$ قاعدة ذات مقدمة أو
  مقدمتين أو ثلاث مقدمات~$!A_1$ و$!A_2$ و$!A_3$، ونتيجتها~$!A$، وكانت
  $\delta_1$ و$\delta_2$ و$\delta_3$~!!{proof}s تنتهي بـ$!A_1$ و
  $!A_2$ و$!A_3$ على الترتيب، فإن كلًا مما يأتي
  \[
  \AxiomC{}
  \RightLabel{$\delta_1$}
  \DeduceC{$!A_1$}
  \RightLabel{$R$}
  \UnaryInfC{$!A$}
  \DisplayProof
\qquad
  \AxiomC{}
  \RightLabel{$\delta_1$}
  \DeduceC{$!A_1$}
  \AxiomC{}
  \RightLabel{$\delta_2$}
  \DeduceC{$!A_2$}
  \RightLabel{$R$}
  \BinaryInfC{$!A$}
  \DisplayProof
\qquad
  \AxiomC{}
  \RightLabel{$\delta_1$}
  \DeduceC{$!A_1$}
  \AxiomC{}
  \RightLabel{$\delta_2$}
  \DeduceC{$!A_2$}
  \AxiomC{}
  \RightLabel{$\delta_3$}
  \DeduceC{$!A_3$}
  \RightLabel{$R$}
  \TrinaryInfC{$!A$}
  \DisplayProof
  \]
  هو~!!a{proof}، على الترتيب، بشرط توافق علامات الفروض في~!!{proof}s،
  أي ألا تحمل~!!{formula}s مختلفة علامة التبرئة نفسها، وأن تكون
  الاستدلالات تطبيقات صحيحة لـ$R$.
  \item تُعد أعلى وقائع~!!{formula} في~!!{proof} الجديد
  \emph{فروضًا مفتوحة} لـ!!{proof} إذا كانت فروضًا مفتوحة في
  $\delta_1$ أو$\delta_2$ أو$\delta_3$، باستثناء الوقائع التي تبرئها
  القاعدة. فإذا وُسمت القاعدة بـ$x$، أو بـ$x\,y$، كانت الوقائع المبرأة:
  في~$\Intro{\lif}$ و$\Intro{\lnot}$ كل الفروض المفتوحة في~$\delta_1$
  الموسومة~$x$؛ وفي~$\Elim{\lexists}$ كل الفروض المفتوحة الموسومة~$x$
  في~$\delta_2$؛ وفي~$\Elim{\lor}$ كل الفروض المفتوحة الموسومة~$x$ في
  $\delta_2$ وكل الفروض المفتوحة الموسومة~$y$ في~$\delta_3$.
\end{enumerate}
\end{defn}

تسمى~!!{proof}s المستخدمة في البناء الاستقرائي لـ!!a{proof}
\emph{!!{proof}s الجزئية} له. وتسمى~!!{proof}s الجزئية المنتهية بمقدمات
استدلال~\emph{!!{proof}s الجزئية المباشرة} للاستدلال. وترد قواعد
\Log{N1c} و\Log{N1i} في~\olref{tab:N1}.

\subfile{rules-N1}

\begin{defn}
يُعرّف \emph{!!{height}}~$\pheight{\delta}$ لـ!!a{proof}~$\delta$
استقرائيًا كما يأتي.
\begin{enumerate}
  \item إذا تألف~$\delta$ من فرض، كان~$\pheight{\delta}=0$.
  \item إذا انتهى~$\delta$ باستدلال ذي مقدمة واحدة و!!{proof} جزئي
  مباشر~$\delta_1$، كان~$\pheight{\delta}=\pheight{\delta_1}+1$.
  \item إذا انتهى~$\delta$ باستدلال ذي مقدمتين و!!{proof}s جزئية
  مباشرة~$\delta_1$ و$\delta_2$، كان~$\pheight{\delta}=
  \max(\pheight{\delta_1},\pheight{\delta_2})+1$.
  \item إذا انتهى~$\delta$ بـ$\Elim{\lor}$ و!!{proof}s جزئية مباشرة
  $\delta_1$ و$\delta_2$ و$\delta_3$، كان~$\pheight{\delta}=
  \max(\pheight{\delta_1},\pheight{\delta_2},\pheight{\delta_3})+1$.
\end{enumerate}
إذا كان للمتتالية~$S$~!!a{proof} ارتفاعه~$\le n$، كتبنا~$\Proves[n] S$.
\end{defn}

\begin{ex}
في~\Log{N1i}، تُعد أي~!!{formula}~!!a{proof} من نفسها بوصفها فرضًا
مفتوحًا. ومن ثم
\[
!A \land !B^x
\]
هو~!!a{proof}~$\delta_1$، وارتفاعه~$0$.
ومن~$\delta_1$ نستطيع الحصول على~!!{proof}s اثنين~$\delta_2$ و
$\delta_3$ بتطبيق إحدى نسختي~$\Elim{\land}$:
\[
\AxiomC{$!A \land !B^x$}
\RightLabel{\Elim{\land}[2]}
\UnaryInfC{$!B$}
\DisplayProof
\qquad
\AxiomC{$!A \land !B^x$}
\RightLabel{\Elim{\land}[1]}
\UnaryInfC{$!A$}
\DisplayProof
\]
والـ!!{proof} الجزئي المباشر للاستدلال الأخير في كل من~$\delta_2$ و
$\delta_3$ هو~$\delta_1$. وفي كل~!!{proof} منهما تكون واقعة
$!A\land !B$ فرضًا مفتوحًا. ولدينا
$\pheight{\delta_2}=\pheight{\delta_3}=1$.

تتطلب قاعدة~$\Intro{\land}$ مقدمتين من الصورتين~$!B$ و$!A$ لتسويغ
نتيجة من الصورة~$!B \land !A$. ولذلك يكون الآتي~!!a{proof}~$\delta_4$:
\[
\AxiomC{$!A \land !B^x$}
\RightLabel{\Elim{\land}[2]}
\UnaryInfC{$!B$}
\LeftSubproofLabel{$\delta_2$}
\AxiomC{$!A \land !B^x$}
\RightLabel{\Elim{\land}[1]}
\UnaryInfC{$!A$}
\RightSubproofLabel{$\delta_3$}
\RightLabel{\Intro{\land}}
\BinaryInfC{$!B \land !A$}
\DisplayProof
\]
وارتفاعه~$\pheight{\delta_4}=\max(\pheight{\delta_2},
\pheight{\delta_3})+1=\max(1,1)+1=2$. وله واقعتان من
$!A \land !B$ بوصفهما فرضين، وكلتاهما مفتوحة.

وأخيرًا نستطيع تطبيق~$\Intro{\lif}$ للحصول على~$\delta_5$:
\[
\AxiomC{$\Discharge{!A \land !B}{x}$}
\RightLabel{\Elim{\land}[2]}
\UnaryInfC{$!B$}
\LeftSubproofLabel{$\delta_2$}
\AxiomC{$\Discharge{!A \land !B}{x}$}
\RightLabel{\Elim{\land}[1]}
\UnaryInfC{$!A$}
\RightSubproofLabel{$\delta_3$}
\RightLabel{\Intro{\land}}
\BinaryInfC{$!B \land !A$}
\DischargeRule{\Intro{\lif}}{x}
\UnaryInfC{$(!A \land !B) \lif (!B \land !A)$}
\DisplayProof
\]
وارتفاعه~$\pheight{\delta_4}+1=3$. ويبرئ استدلال~$\Intro{\lif}$ كل
وقائع الفروض من الصورة~$!A \land !B$ الموسومة~$x$، أي الفرضين اللذين
كانا مفتوحين من قبل في البرهان الجزئي المباشر. ولأنهما الفرضان
الوحيدان لـ$\delta_5$، فلا تكون له فروض مفتوحة، ومن ثم يكون~!!a{proof}
\emph{لـ}!!{formula} نهايته~$(!A \land !B) \lif (!B \land !A)$.
\end{ex}

\end{document}
